\documentclass[10pt,letterpaper]{article}
\usepackage{technical-testing}
\setDocTitle{Walk Runs And Interruptibility Problem Set}
\setDocSubtitle{Original problems with numerical checks and evidence requirements}
\setDocVersion{1.0}
\setDocStatus{Educational / Draft}
\setDocOwner{Game Development}
\setDocDate{2026-10-30}
\setDocSummary{Original problems with numerical checks and evidence requirements}
\begin{document}
\makedoctitle
\makedocsummary

\section{Problems}
\begin{enumerate}
\item Annotate contact, passing, and suspension events on a 24-frame walk. Preserve stationary world-space feet during contact.
\item Construct a looped running clip where its final pose and velocity are compatible with the first sample. Test at two render rates.
\item Define an attack startup, active, recovery and return-to-locomotion lifecycle. Specify interruption guard conditions for jump, dodge, and death.
\item Measure the 95th percentile time between player input and the first readable animation cue over 100 activations. Report hardware, frame pacing, and quantization error.
\end{enumerate}
\section{Acceptance}
Use both animation telemetry and participant perception. A low measured latency does not by itself mean an attack reads correctly; a visual flourish may be attractive but break input responsiveness.
\end{document}
